\section{Context Adherence Metrics}

The evaluation of context adherence in AI assistant responses is crucial for ensuring relevant and appropriate interactions. This section presents a framework for assessing how well an assistant's responses align with the provided context and expected outcomes. Our approach builds upon the LLM-as-a-judge paradigm introduced by Zheng et al.~\cite{zheng2023judging}, which demonstrated that strong LLMs can achieve over 80\% agreement with human preferences---matching inter-human agreement levels---when used as evaluators.

\subsection{The LLM-as-a-Judge Paradigm}

Traditional evaluation of open-ended generative tasks relies on human annotators, which is expensive, slow, and difficult to scale. The LLM-as-a-judge framework~\cite{zheng2023judging} proposes leveraging capable language models as automated evaluators. Zheng et al. showed that models such as GPT-4 can serve as reliable proxies for human judgment across a wide range of evaluation criteria, including relevance, accuracy, and coherence.

Our implementation adopts this paradigm through a dedicated \texttt{Judge} class that wraps any LangChain-compatible language model (\texttt{BaseChatModel}). The judge supports two operational modes:

\begin{enumerate}
    \item \textbf{Structured Output Mode}: Uses the model's native structured output capabilities (via \texttt{create\_agent} with \texttt{response\_format}) to enforce schema validation on judge responses, ensuring that scores and insights always conform to a well-defined Pydantic schema.
    \item \textbf{Regex Extraction Mode}: Injects the expected JSON schema directly into the system prompt, then extracts the structured response from the model's free-form output using regex-based pattern matching between configurable delimiters (\texttt{bos\_json\_clause} and \texttt{eos\_json\_clause}).
\end{enumerate}

This dual-mode design allows the framework to work with models that support constrained decoding (e.g., OpenAI function calling) and those that do not (e.g., open-source models accessed via Ollama or vLLM), providing maximum flexibility across deployment contexts.

\subsection{Context Evaluation Framework}

The context evaluation process employs the judge in a \emph{Contextual Compliance Analyzer} role. The system prompt instructs the judge to evaluate adherence along the following input components:

\begin{enumerate}
    \item \textbf{Context ($C$)}: The provided background information, system instructions, or knowledge that the assistant should adhere to
    \item \textbf{Human Question ($Q$)}: The user's query or input provided to the assistant
    \item \textbf{Assistant Answer ($A$)}: The actual response generated by the AI assistant under evaluation
    \item \textbf{Ground Truth/Observation ($G$)}: The expected or ideal response, used as a reference for comparison
\end{enumerate}

Two distinct evaluation templates are provided depending on whether a ground truth answer or an observation is available:

\begin{itemize}
    \item When a \textbf{ground truth response} is provided, the judge receives both the expected and actual assistant answers, and evaluates whether the actual answer maintains context fidelity compared to the reference.
    \item When an \textbf{observation} is provided instead (e.g., a human annotation or a rubric), the judge incorporates the observation as an additional evaluative criterion, using it to contextualize and support the assessment.
\end{itemize}

\subsection{Structured Output Schema}

The judge is constrained to produce a structured output conforming to the \texttt{ContextJudgeOutput} schema, defined as:

\begin{itemize}
    \item \textbf{Score} ($s \in [0, 1]$): A continuous context alignment score where 0 indicates complete misalignment with the provided context and 1 indicates perfect adherence
    \item \textbf{Insight} ($I$): A free-text explanation of the judge's reasoning, including observations about how well the answer fits the context, potential areas of improvement, or discrepancies noted
\end{itemize}

This schema is enforced either through the model's structured output API or through the JSON extraction pipeline, depending on the configured mode. A retry mechanism (up to 5 attempts) ensures robustness against transient model failures or malformed outputs.

\subsection{Chain-of-Thought Evaluation}

Following the Chain-of-Thought (CoT) framework introduced by Wei et al.~\cite{wei2022chain}, we extract the model's intermediate reasoning process when available. LangChain's \texttt{additional\_kwargs} interface provides access to the \texttt{reasoning\_content} field, which captures the model's step-by-step deliberation before producing the final structured judgment. This reasoning chain includes:

\begin{itemize}
    \item Step-by-step analysis of how each part of the assistant's response relates to the provided context
    \item Systematic comparison between the assistant's answer and the ground truth or observation
    \item Identification of specific context mismatches, hallucinations, or scope exceeding
    \item Justification for the assigned numerical score
\end{itemize}

The explicit extraction of reasoning content serves a dual purpose: it enables human auditors to verify the judge's decision-making process (improving transparency and trust), and it provides actionable feedback for assistant improvement.

\subsection{Session-Level Aggregation}

Context awareness is evaluated at the individual interaction level (per question-answer pair), then aggregated to a session-level score. For a session with $n$ interactions, each producing a score $s_i$, the aggregation process is as follows.

\subsubsection{Weighted Aggregation}

Each interaction $i$ can carry an optional user-defined weight $w_i$. If all weights are provided and sum to 1, they are used directly. If some interactions lack explicit weights, the remaining weight mass is distributed equally among unweighted interactions. If no weights are specified, a uniform distribution $w_i = 1/n$ is applied.

\subsubsection{Statistical Modes}

The framework implements the Strategy Pattern~\cite{gamma1994design} through a pluggable \texttt{StatisticalMode} interface, supporting two aggregation paradigms:

\paragraph{Frequentist Mode.} Returns a point estimate via weighted mean:

\begin{equation}
    \bar{s} = \sum_{i=1}^{n} w_i \cdot s_i
\end{equation}

where $w_i$ are the normalized weights for each interaction.

\paragraph{Bayesian Mode.} Returns a bootstrapped posterior distribution with credible intervals. Given the score vector $\mathbf{s} = (s_1, \ldots, s_n)$ and weight vector $\mathbf{w} = (w_1, \ldots, w_n)$, we draw $S$ Monte Carlo bootstrap samples (default $S = 5000$):

\begin{equation}
    \bar{s}^{(k)} = \frac{1}{n} \sum_{j=1}^{n} s_{I_j^{(k)}}, \quad I_j^{(k)} \sim \text{Categorical}(\mathbf{w}), \quad k = 1, \ldots, S
\end{equation}

The session-level context awareness is then reported with a credible interval:

\begin{equation}
    \hat{s} = \mathbb{E}[\bar{s}^{(k)}], \quad \text{CI}_{1-\alpha} = \left[ Q_{\alpha/2}(\bar{s}^{(k)}), \; Q_{1-\alpha/2}(\bar{s}^{(k)}) \right]
\end{equation}

where $Q_p$ denotes the $p$-th quantile of the bootstrap distribution and $\alpha$ is the significance level (default $\alpha = 0.05$ for 95\% credible intervals).

\subsection{Data Representation}

The evaluation results are represented through a hierarchy of typed data models:

\begin{enumerate}
    \item \textbf{ContextInteraction}: Stores the per-interaction result, identified by a unique \texttt{qa\_id} and its corresponding \texttt{context\_awareness} score $s_i$.
    \item \textbf{ContextMetric}: The session-level aggregate, containing:
    \begin{itemize}
        \item \texttt{session\_id} and \texttt{assistant\_id}: Identifiers linking the metric to the evaluated session and assistant
        \item \texttt{n\_interactions}: The total number of interactions evaluated
        \item \texttt{context\_awareness}: The aggregated session score $\hat{s}$
        \item \texttt{context\_awareness\_ci\_low}, \texttt{context\_awareness\_ci\_high}: Lower and upper bounds of the credible interval (populated only in Bayesian mode; \texttt{None} in Frequentist mode)
        \item \texttt{interactions}: The full list of \texttt{ContextInteraction} records for drill-down analysis
    \end{itemize}
\end{enumerate}

All schemas are implemented as Pydantic~\cite{pydantic2024} models with strict validation (e.g., \texttt{score: float = Field(ge=0, le=1)}), ensuring that scores remain within well-defined bounds and preventing data integrity issues throughout the evaluation pipeline.

\subsection{Bias Mitigation}

The LLM-as-a-judge approach is not without limitations. Zheng et al.~\cite{zheng2023judging} identified several biases in LLM judges, including position bias (preference for the first or second response in pairwise settings), verbosity bias (preference for longer responses), and self-enhancement bias (preference for outputs generated by the same model family). Within our context adherence evaluation, these risks are mitigated by:

\begin{itemize}
    \item Evaluating each response independently (single-answer grading, not pairwise), which eliminates position bias
    \item Constraining the judge to score only context adherence rather than overall quality, reducing the influence of verbosity
    \item Decoupling the judge model from the assistant under evaluation, preventing self-enhancement bias
    \item Providing explicit scoring rubrics within the system prompt to anchor the evaluation criteria
\end{itemize}
